% \section{Variance of Adaptive Rate}

% % In this section, we first analyze the variance of adaptive learning rates, then introduce the empirical results that support our hypothesis -- it is the noise of the adaptive learning rate in the early stage causes the convergence problems.
% In this section, we first analyze the variance of adaptive learning rates, then introduce the empirical results that support our hypothesis -- \emph{Due to the lack of samples in the early stage, the adaptive learning rate has undesirably large variance and can cause convergence issues}.

% To begin with, we first analyze a special case.
% When $t = 1$, we have $\psi(g_1) = \sqrt{\frac{1}{g_i^2}}$. 
% In this paper, we view $\{g_1, \cdots, g_t\}$ as i.i.d. random variables drawn from a Normal distribution $\cN(0, \sigma^2)$.
% % Now, let's move to the special case of $t = 1$. 
% % Since $g_1 \sim \gN(0, \sigma^2)$
% Therefore, $\frac{1}{g_i^2}$ is subject to the scaled inverse chi-squared distribution, $\sics(1, \frac{1}{\sigma^2})$.
% It is easy to find that 
% $\Var[\sqrt{\frac{1}{g_i^2}}] \propto \int_0^{\infty} x^{-1} e^{-x} dx$, which is divergent. 
% It means that, the adaptive ratio can be undesirable large in the first stage of learning.
% % or cause convergence problems. 
% At the same time, setting a small learning rate at the early stage can reduce variance (since $\Var[\alpha x] = \alpha^2\Var[x]$) thus alleviating this problem.
% Therefore, we suggest that the large variance of the adaptive learning rate in the early stage can cause convergence problems.
% % In the following of this section, we first analyze the variance of adaptive learning rates, then introduce the empirical results that supports our hypothesis.

% \subsection{Analysis of Adaptive Learning Rate Variance}
% As mentioned before, Adam uses the exponential moving average to calculate the adaptive learning rate. 
% Assume $\{g_1, \cdots, g_t\}$ are i.i.d., their exponential moving average has a larger variance than their average. 
% Also, in the early stage ($t$ is small), the difference of the exponential weights of $\{g_1, \cdots, g_t\}$ is relative small (up to $1 - \beta_2^{t-1}$).
% Therefore, for ease of analysis, we approximate the moving average as the simple average, \ie, $\psi(g_1, \cdots, g_t) \approx \sqrt{\frac{t}{\sum_{i=1}^t g_i^2}}$ (further analysis on this approximation is conducted in Section~\ref{subsec:sma-ema}). 
% Now, we proceed to its variance analysis. 

% Since $g_i \sim \gN(0, \sigma^2)$, we have $\frac{t}{\sum_{i=1}^t g_i^2} \sim \sics(t, \frac{1}{\sigma^2})$. 
% Therefore, we have:
% \begin{equation}
% \E[\frac{t}{\sum_{i=1}^t g_i^2}] = \frac{t}{(t - 2) \sigma^2} \;(\forall t > 2)
% \quad\mbox{and}\quad
% \Var[\frac{t}{\sum_{i=1}^t g_i^2}] = \frac{2 t^2}{(t - 2)^2 (t - 4) \sigma^4} \;(\forall t > 4).
% \label{eqn:inv-chi-sq}
% \end{equation}

% The analytic form of $\Var[\sqrt{\frac{t}{\sum_{i=1}^t g_i^2}}]$ is derived as Equation~\ref{eqn:analytic-sqrt-var} (in Section~\ref{subsec:approx}).
% Since its calculation is complicated and not numerical stable, we use the following approximation, which is observed to have a high precision in our simulation experiments (discussions are elaborated in Section~\ref{subsec:approx}). 
% Specifically, approximating $\sqrt{x}$ to the first order~\citep{wolter2007taylor}, we have: 
% \begin{equation}
% \sqrt{x} \approx \sqrt{\E[x]} + \frac{1}{2\sqrt{\E[x]}}(x - \E[x])
% \quad\mbox{and}\quad
% \Var[\sqrt{x}] \approx \frac{\Var[x]}{4\E[x]}
% \label{eqn:sqrt-var}
% \end{equation}
% Combining Equation~\ref{eqn:inv-chi-sq} and \ref{eqn:sqrt-var}, we have:
% \begin{equation}
% \Var[\sqrt{\frac{t}{\sum_{i=1}^t g_i^2}}] \approx \frac{t}{2(t - 2) (t - 4) \sigma^2}.
% \label{eqn:variance_appro}
% \end{equation}
% It indicates that, due to the lack of samples (small $t$), in the early stage of training, the variance of the adaptive learning rate is substantially larger than the late stage.
% For example, the variance at $t = 5$ is over $100$ times larger than the variance at $t = 500$ (as in Figure~\ref{fig:var_approx}).
% In the next section, we further design controlled experiments to verify that the large variance is the cause of the convergence problems, and warmup works as a variance reduction technique.

% % \begin{figure}[t]
% % \centering
% % \begin{subfigure}[b]{0.5\textwidth}
% %     \centering
% %     \includegraphics[width=\textwidth]{fig/ppl.pdf}
% %     \caption{\,}
% % \end{subfigure}%
% % \begin{subfigure}[b]{0.5\textwidth}
% %     \centering
% %     \includegraphics[width=\textwidth]{fig/loss.pdf}
% %     \caption{\,}
% % \end{subfigure}
% % \caption{Training of Transformers on the De-En IWSLT'14 dataset. (a) Training perplexity in the first 50k gradient updates. (b) Training loss in the first 50k gradient updates.}
% %     \label{fig:eps2k}
% % \end{figure}

% \subsection{Warmup as Variance Reduction}
% In this section, we design controlled experiments to verify our hypothesis that it is the large variance of the adaptive learning rate causes the convergence issue. 
% Particularly, we design two variants of Adam: Adam-2k and Adam-eps, make comparisons to Adam with warmup and the vanilla Adam (without warmup) on the IWSLT'14 German to English dataset~\citep{cettolo2014report}. 
% In the first two thousand iterations of Adam-2k, only its exponential moving 2nd moments are updated, while its exponential moving average and parameters are fixed\footnote{Different from the previous work~\cite{gotmare2018a}, all parameters and all first moments are frozen during the first two thousand iterations.}; other than this, it follows the original Adam algorithms. 
% To compare with others, its iterations are indexed from -1999 instead of 1.
% As in Figure~\ref{fig:eps2k}, we observe that, after getting these additional two thousand samples to estimate its second moments, Adam avoids the convergence problem met by the vanilla-Adam. 
% Also, comparing Figure~\ref{fig:histogram_2} and \ref{fig:histogram_3}, getting enough samples prevents the gradient distribution from being distorted. 
% These observations verify that, it is the lack of enough samples causes the convergence issue. 

% To further verify our hypothesis, we try to demonstrate that the convergence problem can be avoided by reducing the variance of the adaptive learning rate.
% % To reduce the variance of the adaptive learning rate, a straight forward way is to increase the value of $\epsilon$ in $\hat{\psi}(g_1, \cdots, g_t) = \frac{\sqrt{1-\beta_2^t}}{\epsilon + \sqrt{(1 - \beta_2)\sum_{i = 1}^t \beta_2^{t-i}g_i^2}}$.
% A straight forward way to reduce the variance is to increase the value of $\epsilon$ in $\hat{\psi}(g_1, \cdots, g_t) = \frac{\sqrt{1-\beta_2^t}}{\epsilon + \sqrt{(1 - \beta_2)\sum_{i = 1}^t \beta_2^{t-i}g_i^2}}$.
% Actually, if we assume $\hat{\psi}(.)$ is subject to the uniform distribution, its variance equals to $\frac{1}{12 \epsilon^2}$.
% Therefore, we design an additional variants of the Vanilla-Adam, Adam-eps, which sets $\epsilon$ from a negligible value ($1\times 10^{-8}$) to a non-negligible value ($1\times 10^{-4}$).
% Its performance is also summarized in Figure~\ref{fig:eps2k}. 
% We observe that it avoids the serious convergence problem met by vanilla-Adam and verifies that the convergence problem can be alleviated by reducing the variance of the adaptive learning rate. 
% Besides, similar to Adam-2k, it prevents the gradient distribution distortion (as shown in Figure~\ref{fig:histogram_3}). 
% However, it results in worse performance comparing to Adam-2k or Adam-warmup. 
% We conjecture this is because large $\epsilon$ distorts the original design of Adam, i.e., updates are not uniformly scaled. 
% % As a brief summary, we introduce both theoretical analysis and empirical evidence supports that the adaptive learning rate has undesirably large variance in the first stage of training, and can cause convergence problems. 
% In the next section, we try to fix this issue by rectifying the variance of the adaptive learning rate. 

% \begin{figure}[t]
% \centering
%     \includegraphics[width=\textwidth]{fig/histogram_3.pdf}
% \caption{The histogram of the absolute value of gradients (on a log scale) during the training of Transformers on the De-En IWSLT' 14 dataset. }
% \label{fig:histogram_3}
% \end{figure}
